% !TeX root = ../HSL-Fitness-Tests.tex
% !TeX spellcheck = en_US

\section{Fastest Walk Test}
This fitness test measures the fastest straight autonomous walk from one side of the field to the other.

\subsection{Setup}
This fitness test takes place on one half of a standard SPL field and requires 1 active robot.

The competing robot is placed 50cm behind the starting touchline in line with the penalty mark.
The run ends when the robot crosses the opposite touchline.

\begin{figure}[ht]
    \centering
    \centering
\resizebox{\linewidth}{!}{%
  \FitnessTestFieldM{%
    \draw[orange, line width=4.0pt] (fieldLineTL) -- (fieldLineTR);
    \FitnessTestDot{($(penaltyMarkR) + (0,-\fieldWidth/2 - 0.5)$)}{leaderactive}
  }%
}

    \caption{Fastest Walk Fitness Test (M-Field): the start point is shown in blue and the finishing touchline is highlighted.}
    \label{fig:walk_test_m}
\end{figure}

\subsection{Challenge Execution}
In each attempt, the active robot must walk autonomously in a straight run from the starting side of the field to the opposite touchline.
Code changes are not allowed during an attempt, but are allowed between attempts.

\subsection{Scoring}
The head referee starts the timer once the robot touches the starting touchline and stops the timer once the target touchline is touched.

The lowest time across all attempts is recorded, rounded to the nearest second.

\section{Kick Accuracy Fitness Test}
This fitness test measures the accuracy with which a robot can kick an HSL competition soccer ball toward specified targets.

\subsection{Setup}
This fitness test takes place on a standard SPL field and requires 1 active robot and 1 ball.
The ball is placed on the center of the edge of the goal area.
The competing robot is placed in the center of the goal line facing the ball.

\begin{figure}[ht]
    \centering
    \centering
\resizebox{\linewidth}{!}{%
  \FitnessTestFieldM{%
    \FitnessTestDot{(goalLineR)}{leaderactive}
    \FitnessTestDot{(goalAreaLineR)}{leaderball}
    \FitnessTestDot{(center)}{leadertarget}
    \FitnessTestLabel{($(center) + (-0.45,-0.45)$)}{A}{black}
    \FitnessTestDot{(fieldLineTL)}{leadertarget}
    \FitnessTestLabel{($(fieldLineTL) + (0.55,-0.55)$)}{B}{black}
    \FitnessTestDot{(fieldLineBL)}{leadertarget}
    \FitnessTestLabel{($(fieldLineBL) + (0.55,0.55)$)}{C}{black}
  }%
}

    \caption{Kick Accuracy Fitness Test (M-Field): robot is shown in blue, targets are shown in green, and ball placement is shown in yellow.}
    \label{fig:kick_accuracy_test_m}
\end{figure}

\subsection{Challenge Execution}
Each attempt consists of kicking toward each of the three targets (illustrated in \cref{fig:kick_accuracy_test_m}; S-Field and L-Field variants are provided in \cref{fig:kick_accuracy_test_s,fig:kick_accuracy_test_l}):
\begin{itemize}
    \item The center mark in the center circle
    \item The left corner on the opposite side of the field
    \item The right corner on the opposite side of the field
\end{itemize}
In each attempt, for each kick, participating teams must start the robot behavior of their competing robot via a single button press. The robot then has \qty{\PenaltyShootoutKickTime}{\second} to kick the ball towards the current target.
The competing robot must not touch the ball a second time after the ball has clearly moved, otherwise the attempt ends immediately without scoring.
Once the ball has come to a complete stop, the distance from the ball to the target is measured.

Code changes are allowed when changing targets.

\subsection{Scoring}
After all attempts, the best distance achieved for each target is recorded. 
Accuracy is scored inversely based on distance: the closer the ball stops to the target, the higher the score.
Each target is scored independently, with the lowest distance measurement across attempts recorded for each target.
Distances are rounded to the closest half-centimetre.

\section{Recovery from Fall Fitness Test}
This fitness test measures how quickly a robot can recover and resume normal operation after falling forward or backward.

\subsection{Setup}
This fitness test takes place on a standard SPL field and requires 1 active robot.
The robot is placed on the ground in the center of the field.
Each run starts from a defined ground pose: one run from a front-fall pose and one run from a back-fall pose.

\subsection{Challenge Execution}
Each attempt starts with the robot already on the ground in the assigned pose (flat on the front or back).
The robot must then autonomously recover and stand up.
The test is performed with two separate attempts: one forward fall and one backward fall.

For each fall direction, teams start the robot recovery behavior via a single button press.
Code changes between forward and backward fall attempts are allowed.

\subsection{Scoring}
The head referee starts timing at the button press, and stops the timer when the robot is standing stably on both feet.
The fastest recovery time across all attempts for each fall direction is recorded, rounded to the nearest tenth of a second.
The scores for forward and backward recovery are summed to produce the final fitness test score.

\appendix

\section{Field Size Variants}

\subsection{Fastest Walk Test Variants}
The S-Field and L-Field variants of the Fastest Walk Test are shown in \cref{fig:walk_test_s,fig:walk_test_l}.

\begin{figure}[ht]
    \centering
    \centering
\resizebox{\linewidth}{!}{%
  \FitnessTestFieldS{%
    \draw[orange, line width=4.0pt] (fieldLineTL) -- (fieldLineTR);
    \FitnessTestDot{($(penaltyMarkR) + (0,-\fieldWidth/2 - 0.5)$)}{leaderactive}
  }%
}

    \caption{Fastest Walk Fitness Test (S-Field): the start point is shown in blue and the finishing touchline is highlighted.}
    \label{fig:walk_test_s}
\end{figure}

\begin{figure}[ht]
    \centering
    \centering
\resizebox{\linewidth}{!}{%
  \FitnessTestFieldL{%
    \draw[orange, line width=4.0pt] (fieldLineTL) -- (fieldLineBL);
    \FitnessTestDot{(($(penaltyAreaLT1) + (0.5,0)$)}{leaderactive}
  }%
}

    \caption{Fastest Walk Fitness Test (L-Field): the start point is shown in blue and the finishing touchline is highlighted.}
    \label{fig:walk_test_l}
\end{figure}

\subsection{Kick Accuracy Fitness Test Variants}
The S-Field and L-Field variants of the Kick Accuracy Fitness Test are shown in \cref{fig:kick_accuracy_test_s,fig:kick_accuracy_test_l}.

\begin{figure}[ht]
    \centering
    \input{figs/fitness_tests/kick_accuracy_test_s}
    \caption{Kick Accuracy Fitness Test (S-Field): robot is shown in blue, targets are shown in green, and ball placement is shown in yellow.}
    \label{fig:kick_accuracy_test_s}
\end{figure}

\begin{figure}[ht]
    \centering
    \input{figs/fitness_tests/kick_accuracy_test_l}
    \caption{Kick Accuracy Fitness Test (L-Field): robot is shown in blue, targets are shown in green, and ball placement is shown in yellow.}
    \label{fig:kick_accuracy_test_l}
\end{figure}
